% filepath: /home/kluo/Documents/Repos/mathphy/homework/2026/hw2_sol.tex
% Solutions to Homework 2 — Math Physics Methods (数学物理方法)

\PassOptionsToPackage{quiet}{xeCJK}
\documentclass[11pt]{article}

\usepackage[margin=1in]{geometry}
\usepackage{amsmath,amsthm,amssymb,mathtools, graphicx, multicol, array}
\usepackage{ctex}
\usepackage{fancyhdr}
\usepackage{fontspec}
\usepackage{fourier-orns}
\usepackage{tikz}
\usetikzlibrary{decorations.markings,arrows.meta}

% Chinese punctuation handling (same as hw2.tex)
\let\oldjuhao=。
\catcode`\。=\active
\newcommand{。}{\ifmmode\text{\oldjuhao}\else\oldjuhao\fi}

\let\olddouhao=，
\catcode`\，=\active
\newcommand{，}{\ifmmode\text{\olddouhao}\else\olddouhao\fi}

% Common math symbols
\newcommand{\N}{\mathbb{N}}
\newcommand{\Z}{\mathbb{Z}}
\newcommand{\R}{\mathbb{R}}
\newcommand{\C}{\mathbb{C}}

% Theorem environments
\theoremstyle{remark}
\newtheorem{problem}{}
\theoremstyle{definition}
\newtheorem{solution}{解}

\renewcommand{\headrule}{%
\vspace{-8pt}\hrulefill
\raisebox{-2.0pt}{\quad\decofourleft\decotwo\decofourright\quad}\hrulefill}

\setlength{\headheight}{15.11996pt}
\setlength{\topmargin}{-3.11996pt}

\newcommand{\duedate}{\zhdate{2026/9/29}}
\pagestyle{fancy}

\fancyhead{}
\fancyhead[L]{数学物理方法}
\zhnumsetup{style={Traditional,Financial}}
\fancyhead[R]{课后习题\zhnumber{2} \quad 参考解答}

\fancyfoot{}
\fancyfoot[C]{\thepage}
\fancyfoot[L]{截止日期\duedate}
\fancyfoot[R]{任课教师:罗凯}

\begin{document}
\renewcommand{\labelenumi}{(\arabic{enumi})}
\renewcommand{\labelenumii}{(\arabic{enumi}.\arabic{enumii})}

%==========================================================================
\begin{problem}
证明
\[
|\sinh y| \leq |\sin(x+\imath y)| \leq |\cosh y|.
\]
提示: $\cosh^2 x = 1+\sinh^2 x$.
\end{problem}

\begin{solution}
\textbf{方法 1.}
已知
\[
\sin(x+\imath y)=\sin x\cos(\imath y)+\sin(\imath y)\cos x,
\]
以及
\[
\cosh(\imath z)=\cos z,\qquad \sinh(\imath z)=\imath\sin z,
\]
故 $\cos(\imath y)=\cosh y$，$\sin(\imath y)=\imath\sinh y$，从而
\[
\sin(x+\imath y)=\sin x\cosh y+\imath\sinh y\cos x.
\]
于是
\begin{align*}
\bigl|\sin(x+\imath y)\bigr|
&=\sqrt{\sin^{2}x\cosh^{2}y+\sinh^{2}y\cos^{2}x} \\
&=\sqrt{\sin^{2}x\bigl(1+\sinh^{2}y\bigr)+\sinh^{2}y\cos^{2}x} \\
&=\sqrt{\sinh^{2}y+\sin^{2}x}
\geq|\sinh y|.
\end{align*}
或者
\begin{align*}
\bigl|\sin(x+\imath y)\bigr|
&=\sqrt{\sin^{2}x\cosh^{2}y+(\cosh^{2}y-1)\cos^{2}x} \\
&=\sqrt{\cosh^{2}y-\cos^{2}x}
\leq|\cosh y|.
\end{align*}
综上 $|\sinh y|\leq\bigl|\sin(x+\imath y)\bigr|\leq|\cosh y|$.

\medskip
\textbf{方法 2.}
已知
\[
\sin(x+\imath y)
=\frac{e^{\imath x-y}-e^{-\imath x+y}}{2\imath}
=\frac12\sin x\,(e^{y}+e^{-y})+\frac{\imath}{2}\cos x\,(e^{y}-e^{-y}).
\]
模长为
\[
\bigl|\sin(x+\imath y)\bigr|
=\frac12\sqrt{e^{2y}+e^{-2y}+2(\sin^{2}x-\cos^{2}x)}.
\]
同理
\[
|\sinh y|=\frac12\sqrt{e^{2y}+e^{-2y}-2},\qquad
|\cosh y|=\frac12\sqrt{e^{2y}+e^{-2y}+2}.
\]
因为 $\sin^{2}x-\cos^{2}x=-\cos 2x$，且 $-1\leq\cos 2x\leq 1$，所以
\[
-2\leq 2(\sin^{2}x-\cos^{2}x)\leq 2,
\]
从而
\[
\sqrt{e^{2y}+e^{-2y}-2}
\leq
\sqrt{e^{2y}+e^{-2y}+2(\sin^{2}x-\cos^{2}x)}
\leq
\sqrt{e^{2y}+e^{-2y}+2},
\]
即 $|\sinh y|\leq\bigl|\sin(x+\imath y)\bigr|\leq|\cosh y|$.
\end{solution}

%==========================================================================
\begin{problem}
求解方程
\[
2\cosh^2 z-3\cosh z+1=0.
\]
\end{problem}

\begin{solution}
$2\cosh^{2}z-3\cosh z+1=0$，方程可解得 $\cosh z=1$ 或 $\cosh z=\dfrac12$.

若 $\cosh z=\dfrac{e^{z}+e^{-z}}{2}=1$，可知 $z=2n\pi\imath$，$n\in\mathbb{Z}$.

若 $\cosh z=\dfrac{e^{z}+e^{-z}}{2}=\dfrac12$，令 $z=x+\imath y$，分式展开并化简得
\[
(e^{x}+e^{-x})\cos y+\imath\sin y\,(e^{x}-e^{-x})=1.
\]
上式虚部为 $0$，可知 $\sin y=0$ 或 $e^{x}-e^{-x}=0$.

若 $\sin y=0$，则 $\cos y=\pm 1$，$e^{x}+e^{-x}=\pm 1$，无实数解.

若 $e^{x}-e^{-x}=0$，即 $x=0$，则上式实部化为 $2\cos y=1$，解得 $y=\pm\dfrac{\pi}{3}+2n\pi$.

此时
\[
z=\pm\imath\Bigl(\frac{\pi}{3}+2n\pi\Bigr),\qquad n\in\mathbb{Z}.
\]
\end{solution}

%==========================================================================
\begin{problem}
验证解析函数的实部和虚部满足二维拉普拉斯方程. 若将解析函数 $f(z)$ 写成极坐标的形式
$f(z)=R(r,\theta)\,e^{\imath\Theta(r,\theta)}$，验证
\begin{itemize}
\item $\dfrac{\partial R}{\partial r}=\dfrac{R}{r}\dfrac{\partial\Theta}{\partial\theta}$，
\item $\dfrac{1}{r}\dfrac{\partial R}{\partial\theta}=-R\dfrac{\partial\Theta}{\partial r}$.
\end{itemize}
\end{problem}

\begin{solution}
\textbf{(1)} 令 $f(z)=u(x,y)+\imath v(x,y)$。
由 C-R 方程得
\[
\frac{\partial u}{\partial x}=\frac{\partial v}{\partial y},\qquad
\frac{\partial u}{\partial y}=-\frac{\partial v}{\partial x}.
\]
于是
\[
\frac{\partial^{2}u}{\partial x^{2}}+\frac{\partial^{2}u}{\partial y^{2}}
=\frac{\partial^{2}v}{\partial x\partial y}+\frac{\partial}{\partial y}\Bigl(\frac{\partial u}{\partial y}\Bigr)
=\frac{\partial^{2}v}{\partial x\partial y}-\frac{\partial^{2}v}{\partial x\partial y}=0.
\]
类似有
\[
\frac{\partial^{2}v}{\partial x^{2}}+\frac{\partial^{2}v}{\partial y^{2}}=0.
\]

\medskip
\textbf{(2) 方法 1.}
\[
f(z)=R(r,\theta)\,e^{\imath\Theta(r,\theta)}
=R(r,\theta)\cos\Theta(r,\theta)+\imath R(r,\theta)\sin\Theta(r,\theta).
\]
令
\[
u(r,\theta)=R\cos\Theta,\qquad v(r,\theta)=R\sin\Theta.
\]
极坐标下的 C-R 方程为
\[
\frac{\partial u}{\partial r}=\frac1r\frac{\partial v}{\partial\theta},\qquad
\frac1r\frac{\partial u}{\partial\theta}=-\frac{\partial v}{\partial r}.
\]
将 $u,v$ 分别对 $r$ 和 $\theta$ 求偏导：
\begin{align}
\frac{\partial u}{\partial r}
&=\frac{\partial R}{\partial r}\cos\Theta-R\frac{\partial\Theta}{\partial r}\sin\Theta, \label{eq:ur}\\
\frac{\partial u}{\partial\theta}
&=\frac{\partial R}{\partial\theta}\cos\Theta-R\frac{\partial\Theta}{\partial\theta}\sin\Theta, \label{eq:ut}\\
\frac{\partial v}{\partial r}
&=\frac{\partial R}{\partial r}\sin\Theta+R\frac{\partial\Theta}{\partial r}\cos\Theta, \label{eq:vr}\\
\frac{\partial v}{\partial\theta}
&=\frac{\partial R}{\partial\theta}\sin\Theta+R\frac{\partial\Theta}{\partial\theta}\cos\Theta. \label{eq:vt}
\end{align}
代入 C-R 方程得
\begin{align*}
\Biggl(\frac{\partial R}{\partial r}-\frac{R}{r}\frac{\partial\Theta}{\partial\theta}\Biggr)\cos\Theta
+\Biggl(-R\frac{\partial\Theta}{\partial r}-\frac1r\frac{\partial R}{\partial\theta}\Biggr)\sin\Theta
&=0, \\[4pt]
\Biggl(\frac1r\frac{\partial R}{\partial\theta}+R\frac{\partial\Theta}{\partial r}\Biggr)\cos\Theta
+\Biggl(-\frac{R}{r}\frac{\partial\Theta}{\partial\theta}+\frac{\partial R}{\partial r}\Biggr)\sin\Theta
&=0.
\end{align*}
记
\[
f=\frac{\partial R}{\partial r}-\frac{R}{r}\frac{\partial\Theta}{\partial\theta},\qquad
g=-R\frac{\partial\Theta}{\partial r}-\frac1r\frac{\partial R}{\partial\theta},
\]
则
\[
\begin{cases}
f\cos\Theta+g\sin\Theta=0,\\
-g\cos\Theta+f\sin\Theta=0
\end{cases}
\quad\Longrightarrow\quad
f=0,\quad g=0.
\]
因此
\[
\frac{\partial R}{\partial r}=\frac{R}{r}\frac{\partial\Theta}{\partial\theta},\qquad
\frac1r\frac{\partial R}{\partial\theta}=-R\frac{\partial\Theta}{\partial r}.
\]

\medskip
\textbf{方法 2（用定义证明）。}\\
令 $z=re^{\imath\theta}$，则 $\Delta z=e^{\imath\theta}\,\Delta r+\imath r e^{\imath\theta}\,\Delta\theta$.

沿径向方向：
\begin{align*}
\lim_{\Delta z\to 0}\frac{f(z+\Delta z)-f(z)}{\Delta z}
&=\lim_{\Delta r\to 0}
\frac{R(r+\Delta r,\theta)\,e^{\imath\Theta(r+\Delta r,\theta)}-R(r,\theta)\,e^{\imath\Theta(r,\theta)}}{e^{\imath\theta}\,\Delta r} \\
&=\frac{e^{\imath\Theta(r,\theta)}}{e^{\imath\theta}}\frac{\partial R(r,\theta)}{\partial r}
+\imath\frac{R(r,\theta)\,e^{\imath\Theta(r,\theta)}}{e^{\imath\theta}}\frac{\partial\Theta(r,\theta)}{\partial r}.
\end{align*}

沿角向方向：
\begin{align*}
\lim_{\Delta z\to 0}\frac{f(z+\Delta z)-f(z)}{\Delta z}
&=\lim_{\Delta\theta\to 0}
\frac{R(r,\theta+\Delta\theta)\,e^{\imath\Theta(r,\theta+\Delta\theta)}-R(r,\theta)\,e^{\imath\Theta(r,\theta)}}{\imath r e^{\imath\theta}\,\Delta\theta} \\
&=\frac{e^{\imath\Theta(r,\theta)}}{\imath r e^{\imath\theta}}\frac{\partial R(r,\theta)}{\partial\theta}
+\imath\frac{R(r,\theta)\,e^{\imath\Theta(r,\theta)}}{\imath r e^{\imath\theta}}\frac{\partial\Theta(r,\theta)}{\partial\theta}.
\end{align*}

两式相等得到
\[
\frac{\partial R(r,\theta)}{\partial r}+\imath R(r,\theta)\frac{\partial\Theta(r,\theta)}{\partial r}
=\frac{1}{\imath r}\frac{\partial R(r,\theta)}{\partial\theta}+\frac{R(r,\theta)}{r}\frac{\partial\Theta(r,\theta)}{\partial\theta}.
\]
实部和虚部分别相等得到
\[
\frac{\partial R}{\partial r}=\frac{R}{r}\frac{\partial\Theta}{\partial\theta},\qquad
\frac1r\frac{\partial R}{\partial\theta}=-R\frac{\partial\Theta}{\partial r}.
\]
证毕。
\end{solution}

%==========================================================================
\begin{problem}
若将复变函数 $f(z)=u+\imath v$ 看成是 $x,y$ 的二元函数，再看成是
$z=x+\imath y,\ z^*=x-\imath y$ 的二元函数，试证明 Cauchy--Riemann 方程等价于
\[
\frac{\partial f}{\partial z^{*}}=0.
\]
提示: 利用 $x=(z+z^*)/2,\ y=(z-z^*)/(2\imath)$.
\end{problem}

\begin{solution}
令 $f(z)=u+\imath v$。由
\[
x=\frac{z+z^{*}}{2},\qquad
y=\frac{z-z^{*}}{2\imath}
\quad\Longrightarrow\quad
\frac{\partial x}{\partial z^{*}}=\frac12,\qquad
\frac{\partial y}{\partial z^{*}}=\frac{\imath}{2}.
\]
所以
\begin{align*}
\frac{\partial f}{\partial z^{*}}
&=\frac{\partial f}{\partial x}\frac{\partial x}{\partial z^{*}}
+\frac{\partial f}{\partial y}\frac{\partial y}{\partial z^{*}}
=\frac12\frac{\partial u}{\partial x}+\frac{\imath}{2}\frac{\partial v}{\partial x}
+\frac{\imath}{2}\frac{\partial u}{\partial y}+\frac{\imath^{2}}{2}\frac{\partial v}{\partial y} \\
&=\frac12\Biggl(\frac{\partial u}{\partial x}-\frac{\partial v}{\partial y}\Biggr)
+\frac{\imath}{2}\Biggl(\frac{\partial v}{\partial x}+\frac{\partial u}{\partial y}\Biggr).
\end{align*}
由 C-R 方程，$\dfrac{\partial f}{\partial z^{*}}=0$.
\end{solution}

%==========================================================================
\begin{problem}
证明:若函数 $f(z)$ 在区域 $B$ 内解析，其模为一常数，则函数 $f(z)$ 本身也必为一常数.
\end{problem}

\begin{solution}
\textbf{方法 1.}
由第 3 题，$|f|$ 为常数，则 $R(r,\theta)=R$，$R$ 是一个常数，所以
\[
\frac{\partial R}{\partial r}=\frac{\partial R}{\partial\theta}=0.
\]
由第 3 题的方程可以推出
\[
\frac{\partial\Theta}{\partial\theta}=\frac{\partial\Theta}{\partial r}=0,
\]
所以 $\Theta$ 为一个常数，$f(z)$ 没有对 $(r,\theta)$ 的依赖，即 $f(z)$ 为常数。

\medskip
\textbf{方法 2.}
令 $f(z)=u+\imath v$，因其模为常数，因而有 $u^{2}+v^{2}=C^{2}$。
对 $y$ 求导得到
\[
2u\frac{\partial u}{\partial y}+2v\frac{\partial v}{\partial y}=0.
\]
再对 $y$ 求导得到
\begin{equation}
\frac{\partial u}{\partial y}\frac{\partial u}{\partial y}+u\frac{\partial^{2}u}{\partial y^{2}}
+\frac{\partial v}{\partial y}\frac{\partial v}{\partial y}+v\frac{\partial^{2}v}{\partial y^{2}}=0. \tag{1}
\end{equation}
类似的，对 $x$ 求两次导得到
\begin{equation}
\frac{\partial u}{\partial x}\frac{\partial u}{\partial x}+u\frac{\partial^{2}u}{\partial x^{2}}
+\frac{\partial v}{\partial x}\frac{\partial v}{\partial x}+v\frac{\partial^{2}v}{\partial x^{2}}=0. \tag{2}
\end{equation}
(1)、(2) 相加得
\begin{align*}
&\Biggl(\frac{\partial u}{\partial y}\Biggr)^{2}+\Biggl(\frac{\partial v}{\partial y}\Biggr)^{2}
+\Biggl(\frac{\partial u}{\partial x}\Biggr)^{2}+\Biggl(\frac{\partial v}{\partial x}\Biggr)^{2} \\
&\qquad
+u\Biggl(\frac{\partial^{2}u}{\partial y^{2}}+\frac{\partial^{2}u}{\partial x^{2}}\Biggr)
+v\Biggl(\frac{\partial^{2}v}{\partial y^{2}}+\frac{\partial^{2}v}{\partial x^{2}}\Biggr)
=0.
\end{align*}
由于 $u$、$v$ 均满足 Laplace 方程，所以上式化为
\[
\Biggl(\frac{\partial u}{\partial y}\Biggr)^{2}+\Biggl(\frac{\partial v}{\partial y}\Biggr)^{2}
+\Biggl(\frac{\partial u}{\partial x}\Biggr)^{2}+\Biggl(\frac{\partial v}{\partial x}\Biggr)^{2}=0.
\]
推出
\[
\frac{\partial u}{\partial x}=\frac{\partial u}{\partial y}=\frac{\partial v}{\partial x}=\frac{\partial v}{\partial y}=0,
\]
因此它们都是常数，所以 $f(z)$ 为常数。

\medskip
\textbf{方法 3.}
若 $|f(z)|\equiv C=0$，则明显 $f(z)=0$。
若 $|f(z)|\equiv C\neq 0$，则此时 $f(z)\neq 0$，且 $f(z)\overline{f(z)}\equiv C^{2}$，即
\[
\overline{f(z)}\equiv\frac{C^{2}}{f(z)}
\]
也是解析函数。令 $f(z)=u+\imath v$，则 $\overline{f(z)}=u-\imath v$。
有 $f(z)$ 和 $\overline{f(z)}$ 均在 $B$ 内解析知
\[
u_{x}=v_{y},\quad u_{y}=-v_{x},\qquad
u_{x}=-v_{y},\quad u_{y}=v_{x},
\]
从而
\[
u_{x}=u_{y}=v_{x}=v_{y}=0,
\]
所以 $f(z)$ 为常数。
\end{solution}

%==========================================================================
\begin{problem}
若 $f(z)$ 和 $g(z)$ 在 $z=a$ 点解析，且 $f(a)=g(a)=0$，而 $g'(a)\neq 0$，试证:
\[
\lim_{z\to a}\frac{f(z)}{g(z)}=\frac{f'(a)}{g'(a)}.
\]
\end{problem}

\begin{solution}
因为 $f(z)$、$g(z)$ 在 $z=a$ 点解析，且 $f(a)=g(a)=0$，$g'(a)\neq 0$，所以
\[
\lim_{z\to a}\frac{f(z)}{g(z)}
=\lim_{z\to a}\frac{f(z)-f(a)}{g(z)-g(a)}
=\lim_{z\to a}
\frac{\dfrac{f(z)-f(a)}{z-a}}{\dfrac{g(z)-g(a)}{z-a}}
=\frac{f'(a)}{g'(a)}.
\]
\end{solution}

%==========================================================================
\begin{problem}
找出下列函数的支点，并讨论 $z$ 绕各支点一周回到原处函数值的变化.
\begin{enumerate}
    \item $\sqrt[3]{1-z^3}$,
    \item $\sqrt{\dfrac{z-a}{z-b}}$,
    \item $\ln(z^2+1)$.
\end{enumerate}
\end{problem}

\begin{solution}
\textbf{(1)}
令 $t=1-z^{3}=0$，则
\[
z=e^{2\pi\imath n/3},\qquad n=3k,\ 3k+1,\ 3k+2,\quad k\in\mathbb{Z}.
\]
即三个支点：
\[
\begin{cases}
z_{1}=1, & n=3k,\\
z_{2}=-\dfrac12+\imath\dfrac{\sqrt3}{2}, & n=3k+1,\\
z_{3}=-\dfrac12-\imath\dfrac{\sqrt3}{2}, & n=3k+2.
\end{cases}
\]
绕 $z=z_{1}$ 支点一周，$z=z_{1}+re^{\imath\theta}$，$\theta\to\theta+2\pi$，故 $z-z_{1}=re^{\imath\theta}\to re^{\imath(\theta+2\pi)}$，
\[
f(z)=\sqrt[3]{t}=\sqrt[3]{(z-z_{1})(z-z_{2})(z-z_{3})}
\;\longrightarrow\;
\sqrt[3]{(z-z_{2})(z-z_{3})(z-z_{1})}\,e^{2\pi\imath/3}.
\]

\textbf{(2)}
支点为 $z_{1}=a$，$z_{2}=b$。绕 $z_{1}$，$z-a=r_{1}e^{\imath\theta_{1}}\to r_{1}e^{\imath(\theta_{1}+2\pi)}$，
\[
f(z)=\sqrt{\frac{z-a}{z-b}}
\;\longrightarrow\;
\sqrt{\frac{z-a}{z-b}}\,e^{\imath\pi}.
\]
绕 $z_{2}$，$z-b=r_{2}e^{\imath\theta_{2}}\to r_{2}e^{\imath(\theta_{2}+2\pi)}$，
\[
f(z)\;\longrightarrow\; f(z)\,e^{-\imath\pi}.
\]

\textbf{(3)}
$f(z)=\ln(z+\imath)+\ln(z-\imath)$，支点 $z_{1}=+\imath$，$z_{2}=-\imath$.

绕 $z_{1}=+\imath$ 一周，$z=\imath+r_{1}e^{\imath\theta_{1}}\to z=\imath+r_{1}e^{\imath(\theta_{1}+2\pi)}$，
\[
f(z)=\ln(z+\imath)+\ln(z-\imath)=\ln r_{1}+\imath\theta_{1}+\ln(z+\imath)
\;\longrightarrow\;
\ln r_{1}+\imath\theta_{1}+\ln(z+\imath)+2\pi\imath.
\]
绕 $z_{2}=-\imath$ 一周，$z=-\imath+r_{2}e^{\imath\theta_{2}}\to z=-\imath+r_{2}e^{\imath(\theta_{2}+2\pi)}$，
\[
f(z)=\ln(z+\imath)+\ln(z-\imath)=\ln r_{2}+\imath\theta_{2}+\ln(z-\imath)
\;\longrightarrow\;
\ln r_{2}+\imath\theta_{2}+\ln(z-\imath)+2\pi\imath.
\]

三个小题均有支点 $z=\infty$，绕其一周，函数不变。
\end{solution}

%==========================================================================
\begin{problem}
验证
\[
\oint_C\frac{\mathrm{d}z}{z^{2}+z}=0,
\]
路径 $C$ 由 $|z|=R>1$ 确定的圆.\\
提示: 直接应用柯西定理是错误的，需要分解因式裂项.
\end{problem}

\begin{solution}
\[
\frac{1}{z^{2}+z}=\frac{1}{z(z+1)}=\frac{1}{z}-\frac{1}{z+1},
\]
在 $z=0$、$z=-1$ 处不解析。
\[
\oint_{C}\Biggl(\frac1z-\frac{1}{z+1}\Biggr)\mathrm{d}z
=\Biggl(\oint_{l_{1}}+\oint_{l_{2}}\Biggr)\Biggl(\frac1z-\frac{1}{z+1}\Biggr)\mathrm{d}z
=-\oint_{l_{1}}\frac{\mathrm{d}z}{z+1}+\oint_{l_{2}}\frac{\mathrm{d}z}{z}
=-2\pi\imath+2\pi\imath=0.
\]

\begin{center}
\begin{tikzpicture}[
  scale=1.05,
  circarr/.style={
    decoration={markings, mark=at position 0.13 with {\arrow{stealth}}},
    postaction={decorate}
  }
]
  \draw[-{Stealth},blue!70!black] (-2.7,0) -- (2.85,0);
  \draw[-{Stealth},blue!70!black] (0,-2.35) -- (0,2.5);
  \draw[green!55!black, thick, circarr] (0,0) circle[radius=2.0];
  \node[green!55!black] at (1.55,1.72) {$C$};
  \draw[red!75!black, thick, circarr] (-1,0) circle[radius=0.32];
  \fill[red!75!black] (-1,0) circle[radius=1.3pt];
  \node[red!75!black] at (-1,0.52) {$l_1$};
  \draw[red!75!black, thick, circarr] (0,0) circle[radius=0.32];
  \fill[red!75!black] (0,0) circle[radius=1.3pt];
  \node[red!75!black, anchor=west] at (0.38,0.42) {$l_2$};
  \node[anchor=west] at (2.08,0.28) {$R>1$};
\end{tikzpicture}
\end{center}
\end{solution}

\end{document}
